% GENERATED FILE. Edit canonical lessons, then run npm run generate:books.
% canonical-source-hash: fnv1a64:267f8f7aa7b9bab8
% canonical-lessons: RU-C66-hot, RU-C66-cold, RU-C66-warm, RU-C66-fast, RU-C66-slow

\chapter{Hot, Cold, Warm, Fast, Slow}
\label{ch:ru-hot-cold-warm-fast-slow}
\hlchaptermodality{66}

\begin{chapteropening}
\textbf{By the end of this chapter:} \emph{I can say hot, cold, warm, fast and slow.}

Complete the last lesson of chapter 66: i can say hot, cold, warm, fast and slow.
\end{chapteropening}

\section[goryáchiy]{\ru{горячий} (goryáchiy) — hot}

\label{lesson:RU-C66-hot}



\begin{quote}
Before the new one: say the Russian for long, then the Russian for short.
\end{quote}

\subsection*{You'll want to know: \ru{горячий}}
\textbf{\ru{горячий}} (\emph{goryáchiy}) — \textquotedblleft{}hot\textquotedblright{}.

\begin{grammarlens}[title={What you've built}]
One more word for this chapter.
\end{grammarlens}

\subsection*{Your turn}
\begin{itemize}
\raggedright
  \item \emph{Say it:} \emph{goryáchiy}
  \item \emph{Say it:} \emph{goryáchiy}, once more
  \item \emph{Say it:} say \emph{goryáchiy}
  \item \emph{Recall:} say the Russian for long, then the Russian for short, then say \emph{goryáchiy} again
\end{itemize}

\begin{tcolorbox}[breakable,colback=teal!4,colframe=teal!35!black,title={Before you move on}]
What does \ru{горячий} mean? (\textquotedblleft{}Hot\textquotedblright{}.) Say it once more.
\end{tcolorbox}
\section[kholódnyy]{\ru{холодный} (kholódnyy) — cold}

\label{lesson:RU-C66-cold}



\begin{quote}
Before the new one: say the Russian for short, then the Russian for hot.
\end{quote}

\subsection*{You'll want to know: \ru{холодный}}
\textbf{\ru{холодный}} (\emph{kholódnyy}) — \textquotedblleft{}cold\textquotedblright{}.

\begin{grammarlens}[title={What you've built}]
One more word for this chapter.
\end{grammarlens}

\subsection*{Your turn}
\begin{itemize}
\raggedright
  \item \emph{Say it:} \emph{kholódnyy}
  \item \emph{Say it:} \emph{kholódnyy}, once more
  \item \emph{Say it:} say \emph{kholódnyy}
  \item \emph{Recall:} say the Russian for short, then the Russian for hot, then say \emph{kholódnyy} again
\end{itemize}

\begin{tcolorbox}[breakable,colback=teal!4,colframe=teal!35!black,title={Before you move on}]
What does \ru{холодный} mean? (\textquotedblleft{}Cold\textquotedblright{}.) Say it once more.
\end{tcolorbox}
\section[tyóplyy]{\ru{тёплый} (tyóplyy) — warm}

\label{lesson:RU-C66-warm}



\begin{quote}
Before the new one: say the Russian for hot, then the Russian for cold.
\end{quote}

\subsection*{You'll want to know: \ru{тёплый}}
\textbf{\ru{тёплый}} (\emph{tyóplyy}) — \textquotedblleft{}warm\textquotedblright{}.

\begin{grammarlens}[title={What you've built}]
One more word for this chapter.
\end{grammarlens}

\subsection*{Your turn}
\begin{itemize}
\raggedright
  \item \emph{Say it:} \emph{tyóplyy}
  \item \emph{Say it:} \emph{tyóplyy}, once more
  \item \emph{Say it:} say \emph{tyóplyy}
  \item \emph{Recall:} say the Russian for hot, then the Russian for cold, then say \emph{tyóplyy} again
\end{itemize}

\begin{tcolorbox}[breakable,colback=teal!4,colframe=teal!35!black,title={Before you move on}]
What does \ru{тёплый} mean? (\textquotedblleft{}Warm\textquotedblright{}.) Say it once more.
\end{tcolorbox}
\section[býstryy]{\ru{быстрый} (býstryy) — fast}

\label{lesson:RU-C66-fast}



\begin{quote}
Before the new one: say the Russian for cold, then the Russian for warm.
\end{quote}

\subsection*{You'll want to know: \ru{быстрый}}
\textbf{\ru{быстрый}} (\emph{býstryy}) — \textquotedblleft{}fast\textquotedblright{}.

\begin{grammarlens}[title={What you've built}]
One more word for this chapter.
\end{grammarlens}

\subsection*{Your turn}
\begin{itemize}
\raggedright
  \item \emph{Say it:} \emph{býstryy}
  \item \emph{Say it:} \emph{býstryy}, once more
  \item \emph{Say it:} say \emph{býstryy}
  \item \emph{Recall:} say the Russian for cold, then the Russian for warm, then say \emph{býstryy} again
\end{itemize}

\begin{tcolorbox}[breakable,colback=teal!4,colframe=teal!35!black,title={Before you move on}]
What does \ru{быстрый} mean? (\textquotedblleft{}Fast\textquotedblright{}.) Say it once more.
\end{tcolorbox}
\section[médlennyy]{\ru{медленный} (médlennyy) — slow}

\label{lesson:RU-C66-slow}



\begin{quote}
Before the new one: say the Russian for warm, then the Russian for fast.
\end{quote}

\subsection*{You'll want to know: \ru{медленный}}
\textbf{\ru{медленный}} (\emph{médlennyy}) — \textquotedblleft{}slow\textquotedblright{}. The adjective beside \textbf{\ru{медленно}}, slowly.

\begin{grammarlens}[title={What you've built}]
One more word for this chapter.
\end{grammarlens}

\subsection*{Your turn}
\begin{itemize}
\raggedright
  \item \emph{Say it:} \emph{médlennyy}
  \item \emph{Say it:} \emph{médlennyy}, once more
  \item \emph{Say it:} say \emph{médlennyy}
  \item \emph{Recall:} say the Russian for warm, then the Russian for fast, then say \emph{médlennyy} again
\end{itemize}

\begin{tcolorbox}[breakable,colback=teal!4,colframe=teal!35!black,title={Before you move on}]
What does \ru{медленный} mean? (\textquotedblleft{}Slow\textquotedblright{}.) Say it once more.
\end{tcolorbox}
